\subsection{Stemming and Lemmatisation}
\label{sec:stemming-lemmatization}
\slidesrc{Syntax: RE \& word: 48--65}

\subsubsection{The goal: one form for many}
\slidesrc{Syntax: RE \& word: 48}

Documents use different forms of a word: \emph{organize}, \emph{organizes},
\emph{organizing} (inflectional morphemes, Section~\ref{sec:morphology}); and there are
families of derivationally related words with similar meanings: \emph{democratic},
\emph{democratization} (derivational morphemes). For a search, it is often useful that
a query for one of these words also returns documents containing another word of the
same set.

The goal of both stemming and lemmatisation is to \textbf{reduce inflectional forms, and
sometimes derivationally related forms, of a word to a common base form}:
\begin{align*}
  \text{am, are, is} &\;\to\; \text{be}\\
  \text{car, cars, car's, cars'} &\;\to\; \text{car}\\
  \text{the boy's cars are different colors} &\;\to\; \text{the boy car be differ color}
\end{align*}

\subsubsection{Stemming vs.\ lemmatisation}
\slidesrc{Syntax: RE \& word: 49}

\begin{nbdef}{Stemming and lemmatisation}
\begin{itemize}[nosep]
  \item \textbf{Stemming} is a \emph{crude heuristic} process that chops off the ends of
        words, hoping to get it right most of the time. It often removes derivational
        affixes too. The result, the \emph{stem}, need not be a real word.
  \item \textbf{Lemmatisation} does things properly, with a \emph{vocabulary} and a
        \emph{morphological analysis} of words. It normally removes inflectional endings
        only and returns the base or dictionary form of the word, the \textbf{lemma}.
\end{itemize}
Stemming most commonly collapses derivationally related words; lemmatisation commonly
collapses only the different inflectional forms of a lemma.
\end{nbdef}

\begin{center}
\small
\begin{tabular}{@{}l l l@{}}
\toprule
\textbf{Word} & \textbf{Porter stem} & \textbf{Lemma} \\
\midrule
are          & ar       & be \\
theses       & these    & thesis \\
universal    & univers  & universal \\
university   & univers  & university \\
generalizations & gener & generalization \\
\bottomrule
\end{tabular}
\captionof{table}{A stemmer chops, a lemmatiser looks words up. Stems are not words, and
unrelated words can end up with the same stem.}
\end{center}

\subsubsection{The Porter stemmer}
\slidesrc{Syntax: RE \& word: 50--52}

The \textbf{Porter stemmer} (Martin Porter, 1980) is the classic stemming algorithm for
English. It is a simple \textbf{rule-based} and \textbf{heuristic} method: a cascade of
suffix-rewriting rules such as \textsc{ational} $\to$ \textsc{ate}
(\emph{relational} $\to$ \emph{relate}), applied in sets, in order. The slides present
it in seven steps.

\paragraph{Consonants, vowels and the measure.}
\begin{itemize}[nosep]
  \item A \textbf{consonant} is a letter other than \texttt{a}, \texttt{e}, \texttt{i},
        \texttt{o}, \texttt{u}, and other than \texttt{y} when preceded by a consonant.
        A \textbf{vowel} is any other letter. (So the \texttt{y} of \emph{toy} is a
        consonant, the \texttt{y} of \emph{happy} is a vowel.)
  \item Write $C$ for a sequence of one or more consonants and $V$ for a sequence of one
        or more vowels. Then \emph{every} word has the form
        \[ [C]\,(VC)^{m}\,[V] \]
        where the square brackets mean ``optional''. The number $m$ is the
        \textbf{measure} of the word (or stem): roughly, its number of syllables.
\end{itemize}

\begin{nbexample}{measures}
\begin{center}
\small
\begin{tabular}{@{}l l l c@{}}
\toprule
\textbf{Word} & \textbf{Letters as C/V} & \textbf{Form} & $m$ \\
\midrule
tree     & tr$\cdot$ee                        & $C\,V$                & 0 \\
trouble  & tr$\cdot$ou$\cdot$bl$\cdot$e       & $C\,(VC)\,V$          & 1 \\
troubles & tr$\cdot$ou$\cdot$bl$\cdot$e$\cdot$s & $C\,(VC)(VC)$       & 2 \\
oaten    & oa$\cdot$t$\cdot$e$\cdot$n         & $(VC)(VC)$            & 2 \\
\bottomrule
\end{tabular}
\end{center}
The slides write \emph{troubles} as ``C V CVC'': grouping it as $C\,(VC)(VC)$ gives
$m = 2$.
\end{nbexample}

\paragraph{Rule format.} Each rule has the form \quad(\emph{condition})\; $S_1 \to S_2$
\quad: if the word ends with suffix $S_1$ and the \emph{stem} (what is left after
removing $S_1$) satisfies the condition, $S_1$ is replaced by $S_2$. The conditions are:

\begin{center}
\small
\begin{tabular}{@{}l l@{}}
\toprule
\textbf{Condition} & \textbf{Meaning} \\
\midrule
$m > k$, $m = k$ & the measure of the stem \\
\texttt{*S}  & the stem ends with the letter \texttt{S} (similarly \texttt{*L}, \texttt{*Z}) \\
\texttt{*v*} & the stem contains a vowel \\
\texttt{*d}  & the stem ends with a double consonant (e.g.\ \texttt{-tt}, \texttt{-ss}) \\
\texttt{*o}  & the stem ends with $cvc$, where the second $c$ is not \texttt{w}, \texttt{x}, \texttt{y} \\
\bottomrule
\end{tabular}
\end{center}

The measure condition is what makes the stemmer safe on short words: (\emph{m}$>$0)
\textsc{ness} $\to \varepsilon$ turns \emph{goodness} into \emph{good}, but it would not
touch a hypothetical word whose stem would be empty.

\subsubsection{The seven steps}
\slidesrc{Syntax: RE \& word: 53--58}

In the rules, $\varepsilon$ is the empty string (the suffix is deleted).

\begin{center}
\footnotesize
\begin{tabularx}{\linewidth}{@{}>{\raggedright\arraybackslash}p{1.9cm} >{\raggedright\arraybackslash}X >{\raggedright\arraybackslash}X@{}}
\toprule
\textbf{Step} & \textbf{Rules} & \textbf{Examples} \\
\midrule
1 (plurals) & \textsc{sses}$\to$\textsc{ss};\; \textsc{ies}$\to$\textsc{i};\;
  \textsc{ss}$\to$\textsc{ss};\; \textsc{s}$\to\varepsilon$ &
  caresses$\to$caress, ponies$\to$poni, ties$\to$ti, caress$\to$caress, cats$\to$cat \\
\addlinespace
2a (\emph{-ed}, \emph{-ing}) & ($m{>}0$) \textsc{eed}$\to$\textsc{ee};\;
  (\texttt{*v*}) \textsc{ed}$\to\varepsilon$;\; (\texttt{*v*}) \textsc{ing}$\to\varepsilon$ &
  agreed$\to$agree, feed$\to$feed; plastered$\to$plaster, bled$\to$bled;
  motoring$\to$motor, sing$\to$sing \\
\addlinespace
2b (clean-up, only if the 2nd or 3rd rule of 2a fired) &
  \textsc{at}$\to$\textsc{ate};\; \textsc{bl}$\to$\textsc{ble};\;
  (\texttt{*d} and not \texttt{*L}, \texttt{*S}, \texttt{*Z}) $\to$ single letter;\;
  ($m{=}1$ and \texttt{*o}) $\to$ \textsc{e} &
  conflat(ed)$\to$conflate, troubl(ing)$\to$trouble, hopp(ing)$\to$hop,
  tann(ed)$\to$tan, fall(ing)$\to$fall, fil(ing)$\to$file, fail$\to$fail \\
\addlinespace
3 (\emph{y}) & (\texttt{*v*}) \textsc{y}$\to$\textsc{i} & happy$\to$happi, sky$\to$sky \\
\addlinespace
4 (derivational I) & ($m{>}0$) \textsc{ational}$\to$\textsc{ate};\;
  \textsc{ization}$\to$\textsc{ize};\; \textsc{biliti}$\to$\textsc{ble};\; \ldots &
  relational$\to$relate, generalization$\to$generalize, sensibiliti$\to$sensible \\
\addlinespace
5 (derivational II) & ($m{>}0$) \textsc{icate}$\to$\textsc{ic};\;
  \textsc{ful}$\to\varepsilon$;\; \textsc{ness}$\to\varepsilon$;\; \ldots &
  triplicate$\to$triplic, hopeful$\to$hope, goodness$\to$good \\
\addlinespace
6 (derivational III) & ($m{>}1$) \textsc{ance}, \textsc{ent}, \textsc{ive},
  \textsc{er}, \textsc{al}, \textsc{ant}, \ldots $\to\varepsilon$ &
  allowance$\to$allow, dependent$\to$depend, effective$\to$effect \\
\addlinespace
7a (final \emph{e}) & ($m{>}1$) \textsc{e}$\to\varepsilon$;\;
  ($m{=}1$ and not \texttt{*o}) \textsc{e}$\to\varepsilon$ & probate$\to$probat,
  cease$\to$ceas \\
\addlinespace
7b (double \emph{l}) & ($m{>}1$ and \texttt{*d} and \texttt{*L}) $\to$ single letter &
  controll$\to$control, roll$\to$roll \\
\bottomrule
\end{tabularx}
\captionof{table}{The Porter stemmer as presented in the slides (in Porter's paper the
same rules are numbered 1a, 1b, 1c, 2, 3, 4, 5a, 5b).}
\end{center}

\begin{nbwarning}[corrections to the Porter slides]
\begin{itemize}[nosep]
  \item \textbf{Consonant definition}: the slide reads ``a letter other than A, E, I, O,
        U, and Y preceded by consonant''. Porter's definition is ``other than A, E, I, O,
        U, \emph{and other than} Y preceded by a consonant'': a Y after a consonant is a
        vowel (\emph{happy}), otherwise a consonant (\emph{toy}).
  \item \textbf{Step 2a}: the slide writes the condition of \textsc{eed} as $m>1$, then
        uses \emph{agreed} $\to$ \emph{agree} as an example, but \emph{agr} has $m=1$.
        Porter's rule is $(m>0)$ \textsc{eed} $\to$ \textsc{ee}, as in the table above.
  \item \textbf{Step 7a}: the slide's second rule is ``($m{=}1$ and not \texttt{*o})
        \textsc{ness} $\to \varepsilon$'' with \emph{goodness} as the example (copied
        from step 5). In Porter's algorithm this rule removes a final \textsc{e}:
        \emph{cease} $\to$ \emph{ceas}.
\end{itemize}
\end{nbwarning}

\subsubsection{Examples and problems}
\slidesrc{Syntax: RE \& word: 59--62}

\begin{nbexample}{running the steps}
\begin{itemize}[nosep]
  \item \emph{controlling} $\xrightarrow{\text{2a: } -\textsc{ing}}$ \emph{controll}
        $\xrightarrow{\text{7b}}$ \emph{control}.
  \item \emph{singing} $\xrightarrow{\text{2a: } -\textsc{ing}}$ \emph{sing} (no further
        rule applies).
  \item \emph{computers} $\xrightarrow{\text{1: } -\textsc{s}}$ \emph{computer}
        $\xrightarrow{\text{6: } -\textsc{er}}$ \emph{comput}.
  \item \emph{generalizations} $\xrightarrow{\text{1}}$ \emph{generalization}
        $\xrightarrow{\text{4: \textsc{ization}}\to\textsc{ize}}$ \emph{generalize}
        $\xrightarrow{\text{5: \textsc{alize}}\to\textsc{al}}$ \emph{general}
        $\xrightarrow{\text{6: } -\textsc{al}}$ \emph{gener}.
\end{itemize}
\end{nbexample}

\begin{nbwarning}[the examples slide]
On the slide, \emph{singing} ends with ``Step 6, Rule 4 $\to$ compute'', a line copied
from the \emph{computers} example: \emph{singing} simply becomes \emph{sing}. The slide
also stops \emph{computers} at \emph{compute} and \emph{generalizations} at
\emph{general}. Running the complete algorithm (e.g.\ NLTK's implementation) gives
\emph{comput} and \emph{gener}: stems are not required to be words, and Porter's
algorithm removes \textsc{er} and \textsc{al} in its last derivational step.
\end{nbwarning}

\paragraph{Problems.} Being heuristic, the stemmer makes two kinds of errors:
\begin{itemize}[nosep]
  \item \textbf{over-stemming}: unrelated words get the same stem. \emph{elephants}
        $\to$ \emph{elephant} (step 1) $\to$ \emph{eleph} (step 6 removes \textsc{ant} as
        if it were a suffix); \emph{university} and \emph{universal} both become
        \emph{univers};
  \item \textbf{under-stemming}: forms of the same word get different stems.
        \emph{thesis} $\to$ \emph{thesi} but \emph{theses} $\to$ \emph{these}.
\end{itemize}
The slides also list \emph{doing} $\to$ \emph{doe} as a problem; the standard algorithm
actually returns \emph{do} for \emph{doing} (the \textsc{e} of step 2b is added only when
$m=1$, and \emph{do} has $m=0$).

References: Porter, M.\,F.\ (1980), \emph{An algorithm for suffix stripping},
Program 14(3):130--137; the Porter stemmer home page
\url{https://tartarus.org/martin/PorterStemmer/}; Jurafsky and Martin, Section~3.4.
Online demos: \url{https://9ol.es/porter_js_demo.html}; stemmers for other languages:
\url{http://text-processing.com/demo/stem/}.

\subsubsection{Two texts to test a stemmer}
\slidesrc{Syntax: RE \& word: 63--65}

The slides propose two short texts full of traps. Here is what a real stemmer does with
them.

\begin{nbexample}{Italian: a Snowball stemmer}
\emph{Sono stato in uno stato il cui statuto prevede che la biscia non possa entrare in
una bisca. Una messe di polemiche è piovuta sullo statuto, per le limitazioni messe
alla biscia, che è andata a messa insieme al console per consolarsi.}
\begin{center}
\small
\begin{tabular}{@{}l l l@{}}
\toprule
\textbf{Words} & \textbf{Meanings} & \textbf{Stem} \\
\midrule
stato, stato, statuto & ``been'', ``state'', ``statute'' & \emph{stat} \\
biscia, bisca & ``grass snake'', ``gambling den'' & \emph{bisc} \\
messe, messe, messa & ``harvest'', ``put'' (participle), ``mass'' & \emph{mess} \\
console, consolarsi & ``consul'', ``to console oneself'' & \emph{consol} \\
\bottomrule
\end{tabular}
\end{center}
Every group is over-stemmed: completely unrelated words collapse. Even a lemmatiser
needs context here: the two \emph{stato} have different lemmas (\emph{essere}, the verb;
\emph{stato}, the noun), and only the part of speech tells them apart
(Section~\ref{sec:pos-tagging}).
\end{nbexample}

\begin{nbexample}{English: Miss Marple}
\emph{Once Miss Marple had a cold and she went to the doctor: ``I go!'', she said, and she
read the doctor's doctoral thesis. ``These are good theses!'', and -- although she missed
her train -- Miss Marple understood that University has a universal mission.}

Porter gives: \emph{Miss}, \emph{missed} $\to$ \emph{miss} (a title and a verb merged);
\emph{doctor}, \emph{doctoral} $\to$ \emph{doctor}; \emph{University}, \emph{universal}
$\to$ \emph{univers}; \emph{theses} $\to$ \emph{these}, the same as the pronoun
\emph{These}, while \emph{thesis} $\to$ \emph{thesi}; \emph{has} $\to$ \emph{ha};
\emph{are} $\to$ \emph{ar}. A lemmatiser instead returns \emph{thesis} for
\emph{theses}, \emph{be} for \emph{are}, \emph{have} for \emph{has}, and keeps
\emph{university} and \emph{universal} apart.
\end{nbexample}

\begin{nbintuition}[which one to use]
Stemming is fast, needs no dictionary and increases \emph{recall} in search (more forms
of a word are found), at the price of \emph{precision} (unrelated words are merged).
Lemmatisation is slower and needs a vocabulary and often the part of speech, but its
output is made of real words. Choose according to which kind of error hurts your
application more.
\end{nbintuition}
